\documentclass[11pt,a4paper]{article}
\usepackage[utf8]{inputenc}
\usepackage[T1]{fontenc}
\usepackage{lmodern}
\usepackage{textcomp}
\usepackage{amsmath,amssymb}
\usepackage{graphicx}
\usepackage{booktabs}
\usepackage{longtable}
\usepackage{array}
\usepackage{listings}
\usepackage{xcolor}
\usepackage[margin=1in]{geometry}
\usepackage{float}
\usepackage[hidelinks,breaklinks]{hyperref}
\usepackage{microtype}
\setlength{\parskip}{4pt}
\lstset{basicstyle=\ttfamily\scriptsize,breaklines=true,breakatwhitespace=false,
  postbreak=\mbox{\textcolor{gray}{$\hookrightarrow$}\space},frame=single,
  framerule=0.3pt,rulecolor=\color{gray!60},columns=fullflexible,keepspaces=true,
  showstringspaces=false,numbers=none,xleftmargin=4pt,xrightmargin=4pt,
  literate={\_}{{\_}}1}
\graphicspath{{figures/}}
\providecommand{\tightlist}{}
\title{Baryonic Suppression of the Matter Power Spectrum: Literature Bracket and CRK-HACC Emulator Predictions}
\author{HEP-Genesis-Agent session record \\ \small Project: ramp-research-baryonic-pk}
\date{2026-10-23}
\begin{document}
\maketitle
\noindent\textit{\small This document is an AI-generated account of an interactive session with the HEP-Genesis agent. Its text was written by a language model strictly from the session record; figures and numbers come only from tool results obtained during the session. The supplementary material and the transcript, generated alongside, document exactly what ran. Verify independently before use.}
\vspace{1em}

\begin{abstract}
We report a two-pass research session on baryonic suppression of the matter power spectrum, $S(k)=P_{\rm hydro}/P_{\rm DMO}$. The session had three aims: survey the literature, predict $S(k)$ at $z=0$ with a hydrodynamics emulator for several feedback settings, and assess the consequences for weak-lensing cosmology.

In the first pass no emulator tool was available. That pass therefore built a literature bracket from full-text quotes. Astrid shows no significant suppression at $z=0$. Horizon-AGN shows $<15\%$ suppression at $k\sim10\,h\,{\rm Mpc}^{-1}$. OWLS AGN shows 1\%, 10\% and 30\% suppression at $k\approx0.3$, 1 and $10\,h\,{\rm Mpc}^{-1}$. An X-ray-calibrated BCemu fit gives a maximum suppression of 21--28\% (68\% CL) near $k\sim7\,h\,{\rm Mpc}^{-1}$. FLAMINGO's strongest variants resemble cosmo-OWLS with $\log\Delta T_{\rm AGN}=8.3$--8.5.

In the second pass a CRK-HACC hydrodynamics emulator was found to be connected, and it was evaluated. At a mid-design fiducial AGN point it gives $k(1\%)=0.59$, $k(5\%)=2.11\,h\,{\rm Mpc}^{-1}$ and $S(k=1)=0.979\pm0.001$. $S$ decreases monotonically to $0.791\pm0.006$ at $k=11.27\,h\,{\rm Mpc}^{-1}$, with no minimum inside the emulator range. A strong-AGN corner point adds only $\approx0.4\%$ at $k=1$. A weak-AGN corner point shows a slight enhancement, $S\approx1.011$ near $k\approx2.4\,h\,{\rm Mpc}^{-1}$. Emulator gas fractions in $\sim10^{14}\,h^{-1}M_\odot$ haloes anticorrelate with suppression between the weak and fiducial points, but not between the fiducial and strong points. Suppression weakens and moves to higher $k$ from $z=0$ to $z=1$.

The weak-lensing assessment rests on quoted literature numbers. Reconciling KiDS-1000 with Planck requires $A_{\rm mod}=0.69\pm0.06$; DES Y3 requires $A_{\rm mod}=0.858\pm0.052$. FLAMINGO concludes it is unlikely that baryon physics drives the required modification. No lensing observable or $S_8$ bias was computed in the session. Several emulator-based results were flagged as contested by the session's internal critic, and edge-of-design tests were left incomplete.
\end{abstract}

\section{Introduction}
The user asked the same four-part question twice:
\begin{enumerate}
\item survey recent work on baryonic suppression of the matter power spectrum;
\item use hydrodynamics emulator tools to predict the power-spectrum ratio $S(k)=P_{\rm hydro}/P_{\rm DMO}$ at $z=0$ for a couple of parameter points;
\item explain the AGN-feedback physics responsible for the suppression;
\item assess the impact on weak-lensing cosmology.
\end{enumerate}

In the first pass no emulator was available. That pass produced a literature-only answer and recorded the emulator evaluation as a dead end. In the second pass a CRK-HACC hydrodynamics emulator (server \texttt{serve-hydro.sh}) was found to be connected and exposed a \texttt{Pk-ratio} observable. The emulator part of the question was then answered by direct evaluation.

The emulator was also used for physics cross-checks. These cover gas fraction versus suppression and redshift evolution. A third check, dependence on $\Omega_b/\Omega_m$, turned out not to be supported by the tool. The weak-lensing part was addressed throughout with sourced statements from the literature. No lensing observable was computed in either pass.

This report follows the questions in order and states which items remained open.

\section{Methods}
\paragraph{Literature tools.} All work ran with local dispatch.
\begin{itemize}
\item \textbf{Search.} In the first pass the abstract-search tool (\texttt{search\_arxiv\_abstracts}) returned HTTP 429 (rate-limited) on all three calls. Searches were therefore redone with \texttt{search\_arxiv} (OpenAlex/arXiv API back ends). In the second pass \texttt{search\_arxiv\_abstracts} worked.
\item \textbf{Full texts.} Core papers were downloaded with \texttt{download\_arxiv\_full\_text} and read by character offset with \texttt{read\_arxiv\_chunk}.
\item \textbf{Citations.} Quotes were validated with a \texttt{cite} tool that writes \texttt{references.bib}. Long quotes containing Greek letters, subscripts or ligatures from PDF extraction failed validation; short ASCII fragments passed.
\item \textbf{Figures.} Paper figures were extracted with \texttt{get\_paper\_figures}. That tool stores PNGs and captions only; no tool could view or digitise the images. Numerical values were therefore never read off published curves.
\item \textbf{Test job.} \texttt{run\_test\_job} only draws normal random samples and was not used for physics.
\end{itemize}

\paragraph{Emulator.} In the second pass the \texttt{serve-hydro.sh} server listed 14 observables, including \texttt{Pk-ratio}, \texttt{fGas} and \texttt{HMF}. Its seven-parameter design space is given in Table~\ref{tab:params}. The assistant's board describes the emulator as a SEPIA Gaussian process trained on 110 CRK-HACC hydrodynamical simulations in $400\,h^{-1}{\rm Mpc}$ boxes. That description is the assistant's statement and is not shown in a tool result.

\texttt{Pk-ratio} returns 799 $k$ bins starting at $k=0.0222\,h\,{\rm Mpc}^{-1}$, as mean and $1\sigma$ GP uncertainty. The quantity is labelled $P_{\rm hydro}/P_{\rm grav}$ and is trained at $z=0$, 0.1, 0.5, 1 and 2. \texttt{fGas} returns $M_{\rm gas}/M_{\rm 500c}$ within $R_{\rm 500c}$ in only three $M_{\rm 500c}$ bins: $5.01\times10^{13}$, $7.94\times10^{13}$ and $1.26\times10^{14}\,h^{-1}M_\odot$. The emulator has no $\Omega_b$ parameter.

\begin{table}
\centering
\begin{tabular}{lll}
\toprule
Parameter & Description & Range \\
\midrule
$\kappa_w$ & AGN wind coupling/efficiency & 2--4 \\
$e_w$ & AGN energy efficiency / wind energy fraction & 0.2--1.0 \\
$M_{\rm seed}$ & BH seed mass ($10^6\,M_\odot$) & 0.6--2.0 \\
$v_{\rm kin}$ & kinetic feedback velocity ($10^4$ km/s) & 0.1--1.2 \\
$\epsilon_{\rm kin}$ & kinetic feedback efficiency (units of 10) & 0.02--1.2 \\
$\omega_m=\Omega_m h^2$ & cosmology & 0.12--0.155 \\
$\sigma_8$ & cosmology & 0.7--0.9 \\
\bottomrule
\end{tabular}
\caption{Design space of the HACC hydro emulator as returned by \texttt{describe\_parameters}.}
\label{tab:params}
\end{table}

\paragraph{Parameter points.} All runs fixed $\kappa_w=3$ and $M_{\rm seed}=1.3$. Three AGN settings were used:
\begin{itemize}
\item \textbf{weak}: $e_w=0.2$, $v_{\rm kin}=0.1$ (1000 km/s), $\epsilon_{\rm kin}=0.02$; all at the lower edges of the design;
\item \textbf{fiducial}: $e_w=0.6$, $v_{\rm kin}=0.65$, $\epsilon_{\rm kin}=0.61$; mid-design;
\item \textbf{strong}: $e_w=1.0$, $v_{\rm kin}=1.2$, $\epsilon_{\rm kin}=1.2$; all at the upper edges.
\end{itemize}

Two cosmologies were used. Set A has $\omega_m=0.142$, $\sigma_8=0.81$ and was used for the main $z=0$ curves. Set B has $\omega_m=0.1375$, $\sigma_8=0.80$ and was used for the gas-fraction and $z=1$ cross-checks.

Inside-bound reruns were also generated with Set A cosmology to test edge effects:
\begin{itemize}
\item weak-in: $e_w=0.24$, $v_{\rm kin}=0.155$, $\epsilon_{\rm kin}=0.08$;
\item strong-in: $e_w=0.96$, $v_{\rm kin}=1.145$, $\epsilon_{\rm kin}=1.14$.
\end{itemize}

\paragraph{Reading the output files.} \texttt{load\_table\_data} failed on the output CSVs because of their text header \texttt{x,y\_mean,y\_std}. \texttt{read\_text\_file} truncated them near $k\approx3.7\,h\,{\rm Mpc}^{-1}$. For one file (Set A fiducial) the high-$k$ portion was read via \texttt{read\_arxiv\_chunk} by character offset.

The crossing scales $k(1\%)$ and $k(5\%)$ are the $k$ where $S$ crosses 0.99 and 0.95. They were read off the CSV rows.

One comparison plot was produced with \texttt{plot\_observable\_comparison}. The supplement contains the full tool-call log.

\section{Results}
\subsection{Literature survey and the simulation bracket (first pass)}
Every number in this subsection is a published value quoted from paper text; none was computed in the session.

\paragraph{OWLS AGN.} van Daalen et al.\ (arXiv:1104.1174) state that, at $z=0$, the OWLS AGN run gives ``a decrease in power relative to a dark matter only simulation ranging \ldots\ from 1\% at $k\approx0.3\,h\,{\rm Mpc}^{-1}$ to 10\% at $k\approx1\,h\,{\rm Mpc}^{-1}$ and to 30\% at $k\approx10\,h\,{\rm Mpc}^{-1}$''. Cooling is dominant only for $k\gtrsim70\,h\,{\rm Mpc}^{-1}$. This implies $S(1)\approx0.90$ and $S(10)\approx0.70$. An initial label of OWLS as the ``strong/fiducial'' reference was revised to ``intermediate'' after the critic contested the ordering.

\paragraph{FLAMINGO.} The FLAMINGO response paper (arXiv:2410.17109) gives the following.
\begin{itemize}
\item In the fiducial model ``the position of the minimum of the response moves from $k=6\,h\,{\rm Mpc}^{-1}$ at $z=2$ to $k=10\,h\,{\rm Mpc}^{-1}$ at $z=0$''. The departure from unity begins at $k\approx0.1\,h\,{\rm Mpc}^{-1}$ by $z=0$.
\item For the thermal $f_{\rm gas}$ variants the minimum stays at $k=8$--$10\,h\,{\rm Mpc}^{-1}$. Below $f_{\rm gas}-8\sigma$ the response ``reaches a saturation value''.
\item Collimated jets give a larger response than thermal AGN at fixed $f_{\rm gas}$.
\item The $f_{\rm gas}-8\sigma$ thermal and jet models closely resemble cosmo-OWLS with $\log\Delta T_{\rm AGN}=8.3$ and 8.5. Only Illustris and cosmo-OWLS 8.7 lie beyond the emulator's reach.
\item Both $f_{\rm gas}-8\sigma$ models are ``within the ranges demanded by the analysis of Bigwood et al.\ 2024''. However, they ``display gas fractions that are in strong tension with the gas fractions in groups and clusters inferred from joint X-ray and weak-lensing analyses''.
\item The Amon \& Efstathiou response, and to a lesser extent the Preston et al.\ response, are stated to be out of reach of their emulator.
\end{itemize}
No numerical $S(k=1)$ or $S_{\rm min}$ for the FLAMINGO strong models appears in the text; those amplitudes are only in figures that could not be read.

\paragraph{Weak-end simulations.} Horizon-AGN (arXiv:1801.08559) gives $<15\%$ suppression at $k\sim10\,h\,{\rm Mpc}^{-1}$ at $z=0$. It shows an enhancement by a factor of a few at smaller scales and non-monotonic redshift evolution, and its authors conclude its AGN feedback is probably not strong enough.

Astrid (arXiv:2605.04176) finds ``no significant suppression at $z=0$ and mild suppression at $z=0.2$, weaker than in other state-of-the-art simulations''. At $z=0.2$ the mild suppression extends to $k\sim1\,h\,{\rm Mpc}^{-1}$, and the ratio rises back above unity at $z=0$. Black-hole repositioning ``boosts kinetic AGN feedback (e.g., by a factor of 2 at $z=1.5$)''. At $z=0$ the repositioning run lies $\sim10\%$ below the dynamical-friction run around $k=20\,h\,{\rm Mpc}^{-1}$.

\paragraph{X-ray-calibrated BCemu.} Giri \& Schneider (arXiv:2108.08863, Sec.~4.2) fit the seven-parameter BCemu baryonification model to X-ray gas fractions. The fit assumes Planck cosmology and $b_{\rm hse}=0.26\pm0.07$, with the posterior near $\log_{10}M_c\sim12.6$ and preferring low $\delta$.
\begin{itemize}
\item At 68\% CL, suppression first exceeds 1\% at $k\sim0.1$--$0.4\,h\,{\rm Mpc}^{-1}$ and reaches a maximum of 21--28\% ($S_{\rm min}\sim0.72$--0.79) at $k\sim7\,h\,{\rm Mpc}^{-1}$. The abstract rounds this to 20--28\%.
\item At 99\% CL, $k(1\%)$ lies in 0.1--1.3 and the maximum is 14--38\%.
\item Only cosmo-OWLS 8.5 is fully consistent with the band. cosmo-OWLS 8.0, BAHAMAS and OWLS lie between the 68\% and 99\% regions. IllustrisTNG and Horizon-AGN fall outside the 99\% range because their suppression is shifted to higher $k$.
\end{itemize}






Figure~\ref{fig:21}, reproduced from arXiv:2108.08863 (its Fig.~8), shows this constraint. The left panel has binned X-ray gas and stellar fractions with 68 and 99\% credible regions and simulation curves overlaid. The constraint on $S(k)$ follows from the same MCMC. The figure is the observational anchor for the strong-versus-weak spread: simulations whose gas fractions miss the X-ray data also miss the inferred $S(k)$ band.

The model's ability to reproduce simulation spectra is shown in Figure~\ref{fig:2} (arXiv:2108.08863). Its caption states better-than-1\% agreement at $z<2$, degrading at $k\gtrsim8\,h\,{\rm Mpc}^{-1}$ at $z=2$. Figure~\ref{fig:3}, one panel of arXiv:2108.08863 Fig.~3, illustrates the logic of predicting $S(k)$ from a fit to $f_{\rm gas,500}$ and $f_{\rm star,500}$ alone.











\subsection{AGN-feedback physics (first pass, literature)}
Five papers give the physical picture.
\begin{itemize}
\item \textbf{Gas ejection.} van Daalen et al.\ 2019 (arXiv:1906.00968) analyse 92 OWLS, cosmo-OWLS and BAHAMAS spectra. They attribute large-scale suppression to feedback heating gas and ejecting it to large distances, which is needed to match X-ray groups. A model using only the mean baryon fraction in groups predicts the $z=0$ suppression to within 1\% for $k<1\,h\,{\rm Mpc}^{-1}$.
\item \textbf{Halo-mass dependence.} van Loon \& van Daalen (arXiv:2309.06479) find that AGN feedback ejects matter from inside $r_{\rm 500c}$ to $r_{\rm 200m}$ and possibly beyond in haloes with $M_{\rm 500c}\gtrsim10^{13}\,h^{-1}M_\odot$. Group-sized haloes ($10^{14\pm0.25}\,h^{-1}M_\odot$) dominate for $k\lesssim3\,h\,{\rm Mpc}^{-1}$, and lower-mass haloes take over at smaller scales. A lower AGN heating temperature damps the effect at $k\gtrsim4\,h\,{\rm Mpc}^{-1}$.
\item \textbf{Halo baryon fraction as predictor.} The SP(k) suite (arXiv:2305.09710) shows the dependence on halo baryon fraction across 400 models. In Figure~\ref{fig:23} (its Fig.~2) the $z=0.125$ responses are colour-coded by the baryon fraction of $M_{\rm 200c}=10^{14}\,M_\odot$ haloes. Figure~\ref{fig:24} (its Fig.~4) shows the Spearman correlation between suppression and halo baryon fraction versus $k$. Per its caption, progressively lower-mass haloes correlate best as $k$ increases above $\sim4\,h\,{\rm Mpc}^{-1}$, consistent with arXiv:2309.06479.
\item \textbf{Coupling to cosmology.} Elbers et al.\ (arXiv:2403.12967) report that feedback couples to cosmology through $f_b/c_v^2\sim f_b/(\Omega_m\sigma_8)^{1/4}$.
\item \textbf{Limits of the baryon-fraction predictor.} FLAMINGO (arXiv:2410.17109) reports that SP(k), fed with FLAMINGO's own baryon fractions, predicts a stronger response at $k\approx1$ and a weaker one at $k>3\,h\,{\rm Mpc}^{-1}$. This indicates that baryon fraction alone is not a complete predictor.
\end{itemize}











The internal critic marked all first-pass literature findings derived from abstracts (OWLS, FLAMINGO summaries, van Daalen 2019, van Loon, Horizon-AGN, Elbers, the weak-lensing summaries) as contested. Only the later full-text quote findings (BCemu Sec.~4.2, OWLS verbatim, FLAMINGO verbatim, Astrid, $A_{\rm mod}$ values) were marked supported. The physical narrative above should be read with that status.

\subsection{Emulator predictions at $z=0$ (second pass)}
In the second pass the opening survey text again asserted that no emulator was available. The plan corrected this after \texttt{list\_observables} returned \texttt{Pk-ratio}. The predictions below are actual emulator evaluations.






Figure~\ref{fig:36} overlays the three Set-A curves at $z=0$. Key values are in Table~\ref{tab:hacc}.

On large scales all three points are within $\sim10^{-4}$ of unity. At $k=0.086\,h\,{\rm Mpc}^{-1}$: fiducial $S=0.99986$, strong 0.99981, weak 1.00014. Sub-percent deviations persist to $k\approx0.5$--$0.6\,h\,{\rm Mpc}^{-1}$ for all points.

\paragraph{Fiducial point.} $S$ reaches 0.99 at $0.59$ and 0.95 at $2.11\,h\,{\rm Mpc}^{-1}$. The full CSV read shows a monotonic decline:
\begin{itemize}
\item $S(7.45)=0.8214\pm0.0055$;
\item $S(9.006)=0.8044\pm0.0057$;
\item $S(10.01)=0.7970\pm0.0058$;
\item $S(11.27)=0.7910\pm0.0059$.
\end{itemize}
No minimum or high-$k$ upturn is reached up to at least $k=11.27\,h\,{\rm Mpc}^{-1}$. The last few rows were not read, so the emulator $k_{\max}$ is approximately $11.3\,h\,{\rm Mpc}^{-1}$. $S_{\min}$ and $k_{\min}$ are therefore not defined within the emulator range.

\paragraph{Strong and weak points.} For these, $S(k=5)$ and the high-$k$ behaviour were not read; the files were truncated at $k\approx3.7$.

\begin{table}
\centering
\begin{tabular}{lcccc}
\toprule
Point (Set A) & $k(1\%)$ [$h$/Mpc] & $k(5\%)$ [$h$/Mpc] & $S(k=1)$ & $S(k=3.7)$ \\
\midrule
Weak & --- & --- & $1.003\pm0.003$ & $1.004\pm0.011$ \\
Fiducial & 0.59 & 2.11 & $0.979\pm0.001$ & $0.901\pm0.004$ \\
Strong & 0.53 & 1.86 & $0.975\pm0.002$ & $0.892\pm0.007$ \\
\bottomrule
\end{tabular}
\caption{HACC emulator $S(k)$ at $z=0$ ($\omega_m=0.142$, $\sigma_8=0.81$, $\kappa_w=3$, $M_{\rm seed}=1.3$); uncertainties are $1\sigma$ GP. The weak point never drops below unity and peaks at $\approx1.011\pm0.007$ near $k\approx2.4\,h\,{\rm Mpc}^{-1}$.}
\label{tab:hacc}
\end{table}

The response to the AGN parameters is strongly nonlinear. Almost all the change occurs between the weak and fiducial points. Moving to the strong corner adds only $\approx0.4\%$ at $k=1$, roughly $2\sigma$ of the combined GP uncertainty.

Both the weak-point enhancement and the strong-point saturation occur at corners of the design space. They could be edge or extrapolation artefacts. Inside-bound reruns were generated to test this, but their outputs were never read or compared, so this test was not completed. The critic contested the fiducial, strong/weak and high-$k$ emulator findings.

Comparison with the literature is qualitative only. The fiducial $S(1)=0.979$ is a much smaller suppression than the OWLS AGN quote of 10\% at $k\approx1$. The fiducial $S(10)\approx0.797$ lies inside the BCemu X-ray 68\% band of 21--28\% maximum suppression. The HACC points could not be placed numerically against FLAMINGO, because FLAMINGO's amplitudes appear only in figures.

\subsection{Emulator physics cross-checks (second pass, Set B cosmology)}
\paragraph{Gas fraction versus suppression.} Table~\ref{tab:fgas} lists the emulator gas fractions and $S(k=1.013)$. Going from weak to fiducial, $f_{\rm gas}$ at $7.94\times10^{13}\,h^{-1}M_\odot$ drops by $\approx0.03$ and $1-S(1)$ moves from $\approx-0.002$ to $\approx0.024$. This is qualitatively the halo-baryon-fraction relation of arXiv:1906.00968 and arXiv:2305.09710.

The relation does not hold between fiducial and strong. Their $f_{\rm gas}$ values agree within $1\sigma$; the strong point even has slightly higher $f_{\rm gas}$. Yet the strong point suppresses $\approx0.3\%$ more at $k\approx1$. Possible reasons are that $f_{\rm gas}$ within $R_{\rm 500c}$ in three bins does not capture where the ejected gas resides, or the design-edge concerns noted above. The session did not test either.

\begin{table}
\centering
\begin{tabular}{lcccc}
\toprule
Point (Set B) & $f_{\rm gas}(5.01\times10^{13})$ & $f_{\rm gas}(7.94\times10^{13})$ & $f_{\rm gas}(1.26\times10^{14})$ & $S(k=1.013)$ \\
\midrule
Weak & 0.1337 & $0.1396\pm0.0024$ & 0.1438 & $1.0022\pm0.0025$ \\
Fiducial & 0.0956 & $0.1099\pm0.0011$ & 0.1229 & $0.9766\pm0.0010$ \\
Strong & 0.0985 & $0.1111\pm0.0020$ & 0.1237 & $0.9733\pm0.0018$ \\
\bottomrule
\end{tabular}
\caption{Emulator gas fraction $M_{\rm gas}/M_{\rm 500c}$ within $R_{\rm 500c}$ (masses in $h^{-1}M_\odot$) and $S(k)$ at $z=0$ for $\omega_m=0.1375$, $\sigma_8=0.80$.}
\label{tab:fgas}
\end{table}

\paragraph{Cosmology dependence of the fiducial curve.} The fiducial $S(3.7)$ differs between the two cosmologies: 0.901 in Set A versus 0.892 in Set B. This gives an informal sense of the size of the cosmology dependence. The session did not perform a controlled cosmology scan.

\paragraph{Redshift evolution.} The trends from $z=0$ to $z=1$ (Set B) are summarised below.
\begin{itemize}
\item \textbf{Fiducial.} $S(1.013)$ rises from 0.9766 to $0.9848\pm0.0010$, and $S(3.71)$ from 0.8924 to 0.9458. $k(1\%)$ moves from $\approx0.55$ to $\approx0.78$ and $k(5\%)$ from $\approx1.92$ to $\approx3.09\,h\,{\rm Mpc}^{-1}$.
\item \textbf{Strong.} $S(1.013)$ rises from 0.9733 to $0.9881\pm0.0016$, and $S(3.71)$ from 0.8852 to $0.9620\pm0.0057$. $k(1\%)$ moves from $\approx0.51$ to $\approx0.90$. $k(5\%)\approx1.74$ at $z=0$, while at $z=1$ the curve never reaches 0.95 for $k<3.73$.
\end{itemize}
Suppression therefore weakens and moves to higher $k$ with increasing redshift. This agrees in direction with FLAMINGO's statement that the response strengthens and its minimum moves to smaller scales as $z$ falls (arXiv:2410.17109). At $z=1$ the strong point is less suppressed than the fiducial point at both $k\approx1$ and $k\approx3.7$. This reversal relative to $z=0$ was not investigated further.

These redshift results were marked medium confidence and contested. No plots of the gas-fraction or redshift comparisons were made.

\paragraph{Baryon-fraction dependence.} The emulator has no $\Omega_b$ parameter, so the requested dependence on $\Omega_b/\Omega_m$ could not be tested.

\subsection{Weak-lensing impact}
\paragraph{KiDS-1000.} Amon \& Efstathiou (arXiv:2206.11794) model $P=P_L+A_{\rm mod}(P_{\rm NL}-P_L)$.
\begin{itemize}
\item With the cosmology fixed to Planck TTTEEE $\Lambda$CDM, KiDS-1000 $\xi_\pm$ gives ``a mean $A_{\rm mod}=0.69\pm0.06$'', with $\chi^2_{\min}=265.5$ for 225 data points. The authors note this error neglects the Planck cosmology uncertainty.
\item HMCode2020 instead needs $\Theta_{\rm AGN}=9.17^{+0.2}_{-0.8}$, far outside the BAHAMAS-calibrated range 7.6--8.0.
\item KiDS-1000 shear $S_8=0.759^{+0.024}_{-0.021}$ is about 9\% (2.4--2.7$\sigma$) below Planck $S_8=0.829\pm0.012$. DES Y3 shear gives $S_8=0.772^{+0.018}_{-0.017}$.
\item From their Table~1, KiDS $S_8$ is $0.755\pm0.016$ with feedback switched off and $0.774\pm0.021$ with the fiducial $\Theta_{\rm AGN}$ prior. This shift of $\approx0.02$ ($\approx1\sigma$) is the closest quoted analogue of an $S_8$ bias from ignoring baryons.
\end{itemize}

\paragraph{DES Y3.} Preston, Amon \& Efstathiou (arXiv:2305.09827) find:
\begin{itemize}
\item DES Y3 with a Planck prior and no scale cuts gives $A_{\rm mod}=0.858\pm0.052$, about $2.7\sigma$ from unity;
\item KiDS reanalysed with the same prior gives $0.748\pm0.072$;
\item a naive combination gives $0.820\pm0.042$, i.e.\ suppression required at $\approx4\sigma$;
\item BAHAMAS 8.0 ``is well approximated by $A_{\rm mod}=0.9$'' for $k\lesssim1\,h\,{\rm Mpc}^{-1}$;
\item the required suppression ``is still more extreme than predictions from recent numerical simulations'' and must extend to $k\sim0.2\,h\,{\rm Mpc}^{-1}$.
\end{itemize}

For comparison, the HACC fiducial and strong points give sub-percent suppression at $k\lesssim0.5\,h\,{\rm Mpc}^{-1}$. The session did not convert these into an equivalent $A_{\rm mod}$.

\paragraph{FLAMINGO S8 analysis.} FLAMINGO (arXiv:2309.07959) concludes that baryonic effects are not large enough to remove the $S_8$ tension. It states that some mechanism must slow late-time or non-linear growth ``but that it is unlikely that baryon physics is driving this modification''. Its strongest tension comes from tSZ $y$-$y$, then shear$\times$tSZ, then shear-shear.

\paragraph{Joint and other analyses.} The remaining results come from abstracts and full texts as noted.
\begin{itemize}
\item DES Y3 WL + ACT kSZ (arXiv:2404.06098), using BCEmu, gives $S_8=0.823^{+0.019}_{-0.020}$ and prefers feedback more extreme than most simulations. The fitted baryon fraction at $M_{500}\sim10^{13}\,M_\odot$ is $\approx1.4\sigma$ below HSC-XXL X-ray data. At $4\times10^{13}\,M_\odot$, FLAMINGO, MillenniumTNG and FABLE predict baryon fractions $\approx1.5$, 3.1 and $1.4\sigma$ higher; BAHAMAS agrees best.
\item KiDS-1000$\times$tSZ (arXiv:2109.04458; abstract) prefers strong feedback on its own. Adding shear tightens $S_8$ by 40\% and gives feedback consistent with simulations.
\item DES Y3 harmonic-space shear (arXiv:2203.07128; abstract) gives $S_8=0.793^{+0.038}_{-0.025}$. Extending to $k_{\max}=5\,h\,{\rm Mpc}^{-1}$ constrains feedback while leaving $S_8$ essentially unchanged.
\item An early OWLS-based study (arXiv:1105.1075; abstract) reports that ignoring baryons biases $w_0$ by $\approx40\%$ for a Euclid-like survey.
\end{itemize}
The KiDS$\times$tSZ, harmonic-space and $w_0$ results come from abstracts and were marked medium confidence and contested.

\paragraph{What was not found.} No quoted fractional change in cosmic-shear $C_\ell$ at $\ell\sim500$--5000 was found. The Amon/Preston papers work in $\xi_\pm$ and show the suppression only graphically. No lensing observable was computed from the emulator curves.

\section{Discussion}
The session's central computational result is the HACC $z=0$ suppression curve.

\paragraph{Fiducial shape.} Its shape is broadly similar to the literature picture: onset near $k\sim0.5\,h\,{\rm Mpc}^{-1}$, about 2\% at $k=1$, about 20\% at $k\approx10$. Its amplitude at $k\sim1$ is considerably weaker than OWLS AGN. The amplitude near $k\sim10$ falls within the X-ray-calibrated BCemu 68\% band, which was quoted for $k\sim7$.

Two expected features were not observed. The mild large-scale boost mentioned in the second-pass survey text is not present in the fiducial curve; only the weak point exceeds unity at large scales, by $\sim10^{-4}$. The high-$k$ upturn from cooling is also absent below the emulator's $k_{\max}\approx11.3\,h\,{\rm Mpc}^{-1}$. This is consistent with OWLS placing cooling dominance at $k\gtrsim70\,h\,{\rm Mpc}^{-1}$ and FLAMINGO placing the minimum at $k\approx10$ at $z=0$. The emulator range does not extend far enough to locate a minimum.

\paragraph{Saturation at the strong corner.} The strong-AGN point adds only $\approx0.4\%$ suppression at $k=1$ over the fiducial point, and its cluster gas fraction is not lower. Within this emulator, then, pushing the kinetic and wind

\section*{Data and code availability}
All outputs of this session are collected in the project folder \href{file:///Users/nramachandra/hep-genesis-testing/projects/ramp-research-baryonic-pk}{\texttt{/\allowbreak{}Users/\allowbreak{}nramachandra/\allowbreak{}hep-genesis-testing/\allowbreak{}projects/\allowbreak{}ramp-research-baryonic-pk}}.
The session produced 31 data files; the supplement lists them.
Full documentation of how these results were obtained is in the supplementary material (\texttt{supplement.pdf}: implementation details, the tool-call log, the research board with every finding and its evidence, figures not shown here, and the code index) and the verbatim transcript (\texttt{transcript.pdf}), generated alongside this report.


\end{document}
